\newpage
\phantomsection
\label{prf:supremum-absolute-value-image}
\LRAProofFor{thm:supremum-absolute-value-image}

\begin{remark*}[Return]
\hyperref[thm:supremum-absolute-value-image]{Return to Theorem}
\end{remark*}

\begin{theorem*}[Supremum of the Absolute-Value Image]
Let $A\subseteq\mathbb{R}$ be nonempty and bounded. Then
\[
  \sup |A|=\max\{|\sup A|,|\inf A|\},
  \qquad
  |A|=\{|a|:a\in A\}.
\]
\end{theorem*}

\begin{proof}[Professional Standard Proof]
\LRAProofBodyStart
Let
\[
  s=\sup A,\qquad i=\inf A,\qquad
  M=\max\{|s|,|i|\}.
\]
We prove that $M=\sup |A|$.

First let $a\in A$. Since $i\le a\le s$, if $a\ge 0$, then
\[
  |a|=a\le s\le |s|\le M,
\]
and if $a<0$, then
\[
  |a|=-a\le -i=|i|\le M.
\]
Thus $M$ is an upper bound for $|A|$.

It remains to prove leastness. Since $i\le s$, the number $M$ is either $s$
with $s\ge 0$, or $-i$ with $i\le 0$. Let $\varepsilon>0$.

If $M=s$ and $s\ge 0$, the epsilon characterization of the supremum gives
some $a_\varepsilon\in A$ such that
\[
  s-\varepsilon<a_\varepsilon.
\]
Then
\[
  M-\varepsilon=s-\varepsilon<a_\varepsilon\le |a_\varepsilon|,
\]
and $|a_\varepsilon|\in |A|$.

If $M=-i$ and $i\le 0$, the epsilon characterization of the infimum gives
some $b_\varepsilon\in A$ such that
\[
  b_\varepsilon<i+\varepsilon.
\]
Hence
\[
  M-\varepsilon=-i-\varepsilon<-b_\varepsilon\le |b_\varepsilon|,
\]
and $|b_\varepsilon|\in |A|$.

In every case, for each $\varepsilon>0$ there is an element of $|A|$ greater
than $M-\varepsilon$. Together with the upper-bound property above, the epsilon
characterization of the supremum yields
\[
  \sup |A|=M=\max\{|\sup A|,|\inf A|\}.
\]
\end{proof}

\begin{proof}[Detailed Learning Proof]
\LRAProofBodyStart
Set
\[
  s=\sup A
  \qquad\text{and}\qquad
  i=\inf A.
\]
Because $A$ is nonempty and bounded, both $s$ and $i$ exist, and
\[
  i\le a\le s
  \qquad\text{for every }a\in A.
\]
We claim that
\[
  M=\max\{|s|,|i|\}
\]
is the supremum of the absolute-value image
\[
  |A|=\{|a|:a\in A\}.
\]

First we show that $M$ is an upper bound for $|A|$. Choose any element of
$|A|$. It has the form $|a|$ for some $a\in A$. If $a$ is nonnegative, then
$|a|=a$, and the inequality $a\le s$ gives
\[
  |a|=a\le s\le |s|\le M.
\]
If $a$ is negative, then $|a|=-a$. Since $i\le a$, multiplying by $-1$
reverses the inequality and gives $-a\le -i$. In this case $i<0$, so
$-i=|i|$, and therefore
\[
  |a|=-a\le -i=|i|\le M.
\]
Thus every element of $|A|$ is at most $M$.

Now we prove that this upper bound is least. The largest absolute value at the
endpoints must come from one of two endpoint directions: either the right
endpoint contributes $M=s$ with $s\ge 0$, or the left endpoint contributes
$M=-i$ with $i\le 0$. These two alternatives cover all possibilities because
$i\le s$.

Let $\varepsilon>0$. Suppose first that $M=s$ and $s\ge 0$. Since
$s=\sup A$, the epsilon characterization of the supremum gives an element
$a_\varepsilon\in A$ with
\[
  s-\varepsilon<a_\varepsilon.
\]
Taking absolute values can only make a real number larger than or equal to
itself, so
\[
  M-\varepsilon=s-\varepsilon<a_\varepsilon\le |a_\varepsilon|.
\]
The number $|a_\varepsilon|$ belongs to $|A|$, so $|A|$ reaches above
$M-\varepsilon$.

Suppose instead that $M=-i$ and $i\le 0$. Since $i=\inf A$, the epsilon
characterization of the infimum gives an element $b_\varepsilon\in A$ with
\[
  b_\varepsilon<i+\varepsilon.
\]
Multiplying by $-1$ reverses this inequality:
\[
  -i-\varepsilon<-b_\varepsilon.
\]
Since $-b_\varepsilon\le |b_\varepsilon|$, we get
\[
  M-\varepsilon=-i-\varepsilon<-b_\varepsilon\le |b_\varepsilon|.
\]
Again $|b_\varepsilon|\in |A|$, so $|A|$ reaches above $M-\varepsilon$.

We have shown two things: $M$ is an upper bound for $|A|$, and for every
$\varepsilon>0$ the set $|A|$ contains an element greater than
$M-\varepsilon$. By the epsilon characterization of the supremum,
\[
  \sup |A|=M=\max\{|\sup A|,|\inf A|\}.
\]
\end{proof}

\begin{remark*}[Proof structure]
The upper-bound half uses $i\le a\le s$ and splits by the sign of $a$. The
leastness half uses the epsilon characterization at whichever endpoint has
larger absolute value: near $\sup A$ if the right endpoint controls the
maximum, and near $\inf A$ if the left endpoint controls it.
\end{remark*}

\begin{dependencies}
\hyperref[def:absolute-value]{Absolute Value};
\hyperref[def:supremum]{Supremum};
\hyperref[def:infimum]{Infimum};
\hyperref[def:maximum]{Maximum};
\hyperref[def:epsilon-characterization-of-supremum]{Epsilon Characterization of Supremum};
\hyperref[def:epsilon-characterization-of-infimum]{Epsilon Characterization of Infimum}.
\end{dependencies}

\clearpage
